% Source: https://artofproblemsolving.com/community/c3h1607829_difficult_questions_from_nc_state_math_contest
% Stable identity: aops-topic-1607829-q69
\textbf{Range of $f$.} Put $t=2^x$, which takes every positive value except $1$. Then $f=\dfrac{t+1}{t-1}=1+\dfrac{2}{t-1}$.
\begin{itemize}
\item For $t>1$ (i.e.\ $x>0$), $t-1$ runs over $(0,\infty)$, so $f$ runs over $(1,\infty)$.
\item For $0<t<1$ (i.e.\ $x<0$), $t-1$ runs over $(-1,0)$, so $\dfrac{2}{t-1}$ runs over $(-\infty,-2)$ and $f$ over $(-\infty,-1)$.
\end{itemize}
So $f(x)=u$ takes exactly the values $u\in(-\infty,-1)\cup(1,\infty)$; in particular $u\ne0$, so $f(u)$ is defined.

\textbf{Range of $f(u)$ for these $u$.} With $s=2^u$:
\begin{itemize}
\item $u>1$: $s\in(2,\infty)$, $s-1\in(1,\infty)$, $f(u)=1+\dfrac{2}{s-1}\in(1,3)$.
\item $u<-1$: $s\in\left(0,\dfrac{1}{2}\right)$, $s-1\in\left(-1,-\dfrac{1}{2}\right)$, $\dfrac{2}{s-1}\in(-4,-2)$, so $f(u)\in(-3,-1)$.
\end{itemize}
Each of these is a continuous, strictly monotone function of $u$ on an interval, so all values in between occur. The range of $f(f(x))$ is
\[
\boxed{(-3,-1)\cup(1,3)} .
\]
(For example $f(f(1))=f(3)=\dfrac{9}{7}$ lies in $(1,3)$; this already shows that the answer $\left(-\dfrac{1}{3},\dfrac{1}{3}\right)$ quoted in the post is wrong.)
